\section{Experiments}
\label{sec:experiments}

\subsection{Experimental Setup}
\label{sec:setup}

We evaluate seven LLMs across three families---Llama~3.2~3B, Llama~3.1~8B, Llama~3.3~70B~\cite{meta2024llama3}, Mistral~7B~\cite{jiang2023mistral}, Qwen~2.5~7B/14B~\cite{qwen2025qwen25}, Claude Haiku~4.5~\cite{anthropic2025haiku45}---at $n=50$--100 balanced trials per condition, plus Qwen~2.5~32B in the scale sweep and EXP-R1. Claude Haiku~4.5 is the interrogator and feature extractor throughout (cross-family for every target but itself). Claims are 50 matched true/false factual pairs (``Water boils at 100\textdegree C'' vs.\ ``at 85\textdegree C''; Appendix~\ref{app:claims}), tested \emph{instructed} and \emph{prompt-equalized}. Model diversity here shows the failure is not tied to one detector/model pair; it is \emph{not} an exhaustive survey of frontier models. \textbf{Four experiments carry the argument}---§\ref{sec:r1_faithful}, §\ref{sec:prompt_equalized}, §\ref{sec:hedging_baseline} and §\ref{sec:external_audit}---with the factorial of §\ref{sec:factorial} a \emph{diagnostic} that sizes what they leave unattributable; Table~\ref{tab:what_establishes} states what each establishes and the nearest overclaim it does not, Table~\ref{tab:experiment_summary} the full inventory.

\begin{table}[t]
\centering
\caption{What each experiment establishes---and the nearest conclusion it does \emph{not} support. Four experiments carry the argument (§\ref{sec:r1_faithful}, §\ref{sec:prompt_equalized}, §\ref{sec:hedging_baseline}, §\ref{sec:external_audit}); the factorial is a diagnostic and transfer and human validation are in §\ref{sec:discussion}. Results are given in the text. \emph{Unattributable} means the design cannot assign an effect to deception rather than to compliance with the instruction that produced it (§\ref{sec:identification}).}
\label{tab:what_establishes}
\begin{scriptsize}
\setlength{\tabcolsep}{4pt}
\begin{tabular}{@{}p{1.35in}p{3.90in}@{}}
\toprule
\textbf{Experiment} & \textbf{What it establishes---and what it does not} \\
\midrule
\textbf{1--2.} Both detectors collapse under equalization (EXP-R1, R1b/R1c; EXP-A) & The collapse belongs to the \emph{paradigm}, not one pipeline: two mechanistically unrelated detectors fail under one control, on two claim sets and two probe banks. \emph{Not}: that no deception-specific signal exists, nor that the residual \emph{is} compliance \\
\addlinespace[1pt]
\emph{Diagnostic:} $2\!\times\!2$ factorial & The size of the unattributable component, and that output moves within one identically-labelled condition. \emph{Not}: that latent $D$ was held fixed, nor a causal attribution \\
\addlinespace[1pt]
\textbf{3.} Regex baseline, adversary, and belief strata (EXP-B, pre-registered; \emph{observational}) & Measured accuracy is achievable from a surface feature that is neither deception-specific nor paraphrase-robust, and part of the residual tracks whether the model \emph{knows} the claim, where no deception exists. \emph{Not}: that no latent deception signal exists; not a causal $B\!\to\!S$ effect \\
\addlinespace[1pt]
\textbf{4.} External audit of 10 published rollout sets (EXP-XA) & Criterion 4 gets a verdict from data we did not collect, and the channel scored is itself an audit criterion. \emph{Not}: that these corpora hold no deception-specific signal---only that this rule finds none in the channel a deployed detector reads \\
\bottomrule
\end{tabular}
\end{scriptsize}
\end{table}

\paragraph{Statistical analysis policy.} \textbf{Confirmatory} (permutation-tested): the equalization collapse (EXP-R1, EXP-A), below-chance factorial detection, cross-model LOO. \textbf{Pre-registered}: EXP-R1b, EXP-B, EXP-XA's confirmatory arm. All else is exploratory or descriptive. The inferential unit is the \textbf{target~$\times$~claim-pair}: trials within a claim pair are clustered and never treated as independent, so every cross-validation split is grouped by claim pair and permutations are drawn within it. Pooled accuracies are descriptive summaries, not tests.

\subsection{Experiment 1: The Prior Detector Collapses Under Equalization (EXP-R1)}
\label{sec:r1_faithful}

\citet{pacchiardi2023catch} score no behavioral features and never read the target's answer about the claim: after the target defends a statement they ask a \emph{fixed battery of unrelated yes/no elicitation questions}, binarise each answer, and fit one logistic regression. We reproduce this with an order-stable bank of 16 probes across their three published categories (5 lie-related, 6 factual, 5 ambiguous) and no LLM feature scoring; their original bank was unavailable, so this is a \textbf{mechanism-faithful reimplementation}, not a per-question replication (Appendix~\ref{app:pacchiardi}). \textbf{Stratified 5-fold} accuracy is primary (LOO is biased \emph{below} chance in the no-signal case tested here); significance is a 1000-draw label-permutation test (Appendix~\ref{app:r1_protocol}).

\paragraph{Result.} \textbf{The selection rule is fixed by the instructed cell, before any equalized result is seen}: a target on which the mechanism never worked cannot exhibit a collapse. Instructed, it works on five of six targets at 70.0--100.0\% (all $p\!\leq\!0.002$; Table~\ref{tab:r1_faithful}), bracketing the 67--73\% of \citet{pacchiardi2023catch}, so the mechanism functions in its native regime. \emph{Qwen~2.5~7B} is the sixth: the detector never separates its conditions (54.0\%, $p\!=\!0.299$) and equalized all 50 trials return the \emph{identical} 16-bit vector---a degenerate probe channel, not an accuracy of 0\%. Among the remaining five, equalization drops four to chance (48.0--52.0\%, $p\!\geq\!0.425$); pooled, 81.7\%$\to$53.7\%. The fifth, \emph{Qwen~2.5~14B}, holds at 70.0\%, which the pre-registered replication below does not reproduce.

\paragraph{The Qwen~2.5~14B exception does not replicate (EXP-R1b).} Of the 16 dimensions only dim~0, the confidence probe, varies under equalization, and it alone yields the full 70.0\% (YES on 76\% of true- vs.\ 36\% of false-claim trials). Both cell and dimension were selected post hoc, so we replicated with the analysis fixed in advance. \textbf{Pre-registered, one-sided, $\alpha\!=\!0.05$:} \emph{Qwen~2.5~14B retains above-chance equalized accuracy on a disjoint claim set and an independently worded probe bank, with the surviving variance again concentrating in the confidence probe}---with Qwen~14B instructed and Mistral~7B equalized as positive and negative controls.

\textbf{The hypothesis is refuted.} On 50 fresh claim pairs and 16 re-worded probes the mechanism still reaches \textbf{100.0\%} instructed ($p\!=\!0.001$), so the instrument is live; equalized it falls to \textbf{46.0\%} ($p\!=\!1.000$), with 1 of 16 dimensions non-constant and the re-worded confidence probe carrying \emph{no} separation. Extending the same v2 materials to all six targets (\textbf{EXP-R1c}, pre-registered before any new cell was collected) completes the sweep: all six are informative under v2 (86.0--100.0\%, every $p\!=\!0.001$) and all six fall to chance equalized (24.0--58.0\%, every $p\!\geq\!0.173$; pooled 97.0\%$\to$43.3\%), so \textbf{equalization collapses the detector on every target where it had instructed signal, on two independently constructed claim sets and two independently worded probe banks} (Appendix~\ref{app:r1b}). The confound belongs to the \emph{paradigm}, not one implementation: two mechanistically unrelated detectors---prior work's binary-probe mechanism and our pipeline (30--41\,pp; §\ref{sec:prompt_equalized})---fail under the same control. \textbf{One reading of any such collapse is that equalization removes the deception, not the confound}; no manipulation of $E$ can rule that out, which is why the factorial only sizes what is at stake and criterion~4 is where it would be settled.

\subsection{Experiment 2: An Independent Detector Collapses Under the Same Control}
\label{sec:prompt_equalized}
\label{sec:instructed_matched}

This is a \emph{second instance} of §\ref{sec:r1_faithful}'s result, not a separate finding. Instructed, LOO-recalibrated accuracy reaches 93.9--100\%; under EXP-A's identical neutral prompts, with only claim truthfulness varying, accuracy falls to 61\% (3B)--84\% (70B), all five features collapsing (Figure~\ref{fig:feature_collapse}, Appendices~\ref{app:instructed_results},~\ref{app:scale_details}), and matched claim pairs put the instruction-attributable difference at $+$7.5--15\,pp at $\geq$14B and $+$26--31\,pp at $\leq$8B (Appendix~\ref{app:expg_details}). Two controls bound this: withholding an \emph{entire} claim pair moves accuracy 0.3\,pp (\textbf{67.4\%}, $n\!=\!639$), so the signal is not claim-pair memorization; and cross-family re-extraction (criterion 2) gives \textbf{64.7}/66.9\% against Haiku's 74.3\%, with five self-family controls within 0--3\,pp---an extractor-quality gap, not same-family bias (Appendix~\ref{app:supporting_diagnostics}).

\subsection{Diagnostic: Sizing the Unattributable Component}
\label{sec:factorial}

The 2$\times$2 factorial addresses that second reading by crossing truth-value with the deception instruction ($n\!=\!50$/cell, 3 models: \textbf{T/H}, \textbf{F/H}, \textbf{T/D} true claim argued against, \textbf{F/D} standard instructed-lie). It cannot adjudicate (a) against (b)---complying means asserting what the target does not hold true, so $D$ follows $E$ in \emph{both} instructed cells---but it sizes what is at stake, which is why it is a diagnostic here and not one of the four experiments.

\paragraph{Result 1: the elicitation drives a large, consistently signed component of the response.} We regress detector output on $V$, $E$ and their interaction (OLS; two standardized outcomes, the R1-faithful decision score under claim-held-out calibration and the raw marker count). On the complete detector the instruction effect is large, consistently signed, and significant on all three targets ($\beta_E\!=\!+1.27$, $+1.56$, $+1.38$\,SD; all $p\!<\!0.001$), while the response to truth-value is smaller and \emph{changes sign} ($-0.40$, $-0.96$, $+0.07$\,SD). By the Proposition no value of $\beta_E$ distinguishes the two paths, so this sizes the unattributable component.

\paragraph{Result 2: within one instructed-deception condition, both detectors still move.}
\label{sec:td_vs_fd}
\textbf{T/D} and \textbf{F/D} carry the same deception instruction, complied with in both, so the benchmark assigns them the \emph{same} label and only truth-value differs. Latent deception is not observed, so this fixes the \emph{condition}, not $D$ itself: compliance is graded, as the corrective leakage below shows. The marker nonetheless fires on 62\% of Qwen~14B's T/D trials against 6\% of its F/D trials ($p\!=\!0.0001$), and the detector's score shifts $-1.37$\,SD (Mistral) and $+1.61$\,SD (Qwen), both $p\!=\!0.0001$---as large as the instruction main effect and in \emph{opposite} directions. The mechanism is visible: told to argue against a true claim the target corrects it anyway, firing the patterns the rule keys on, so benchmark-calibrated accuracy here is 39.0/40.5/44.0\%---below chance. The contrast confounds truth-value with compliance difficulty and identifies neither, but on no reading is the instruction the ground truth it is taken for (Appendix~\ref{app:factorial_txd}).

\subsection{Experiment 3: What Survives Is Surface-Lexical and Knowledge-Related}
\label{sec:hedging_baseline}

\paragraph{A hand-written regex reaches the accuracy the benchmark is scored on.} A single regex feature---refusal/correction marker count from 14 patterns (``I cannot,'' ``incorrect'')---reaches \textbf{79.8\%} average under equalization with \emph{no learned parameters} ($k\!=\!1$ follows from a prior, not fitting, so it cannot leak across claims), against the prior detector's 81.7\% instructed and 53.7\% equalized, and matching or exceeding our pipeline on 6 of 7 models. The figure depends on the pattern inventory---dropping the most load-bearing pattern gives 69.1\%---so we report \textbf{69--80\% depending on the inventory} and treat the rule as a methodological instrument, not a production detector (Appendices~\ref{app:regex_patterns},~\ref{app:validation_tests}).

\paragraph{The channel is lexical, and does not survive an adversary or rewording.}
\label{sec:adversarial_paraphrase}\label{sec:adversarial_fullregex}\label{sec:paraphrase_invariance}
A style-constraint adversary at two strengths (Llama~3B, Mistral~7B, Qwen~14B; Table~\ref{tab:adv_multimodel}) has a mechanical effect: on Llama~3.2~3B the fraction of \emph{lie} trials carrying a correction marker falls 84\%$\to$48\%$\to$16\% (Appendix~\ref{app:q1_refusal_dist}), with Mistral~7B a near-parity exception (Appendix~\ref{app:mistral_adv_mech}). What is lost is \emph{lexical} signal, not the behavior: across 30 hand-coded adversarial lie trials all 30 still correct the false claim (Appendix~\ref{app:adv_qualitative}), so a semantic detector would be unaffected. Across 12 deception and 5 honest instruction paraphrases the rule sits below chance on nearly every variant (35--51\%): what is tracked is the instruction---its presence or its wording. \textbf{Measured accuracy is thus achievable from a surface feature that is neither deception-specific nor paraphrase-robust}, removable by an adversary who never has to become a better liar.

\paragraph{A substantial component is associated with knowledge conflict (EXP-B).}
\label{sec:belief_strata}
EXP-B holds the neutral prompt fixed---so no cell contains a deception instruction---and varies belief, with strata assigned \emph{per target} by direct query. Pre-registered, the rule reaches \textbf{70.0\%} [62.2--76.8\%] on KNOWN pairs against 57.6\% [48.6--66.2\%] on UNKNOWN: $+$12.4\,pp ($p\!=\!0.0045$, $n\!=\!268$, positive on all three targets), surviving length matching. This is a \emph{stratified observational} contrast on behavioral belief estimates, not $\mathrm{do}(B)$---the strata also differ in topic and surface plausibility---so rather than a causal $B\!\to\!S$ effect it associates part of the residual with whether the model \emph{knows} the claim, in a design containing no deception at all (Appendix~\ref{app:belief_strata}).

\subsection{Experiment 4: No Published Benchmark Supplies Criterion 4 (EXP-XA)}
\label{sec:external_audit}

\begin{table}[t]
\centering
\caption{EXP-XA: criterion 4 on ten rollout sets released with \citet{goldowskydill2025probes}, scored by §\ref{sec:hedging_baseline}'s rule offline with no model calls. Gates are pre-specified: $E$ counts as fixed only if every rollout shares one byte-identical system prompt, and the rule must fire on 5--95\% of some cell of at least 10 rollouts to be applicable, so a rule undefined outside its native format is reported as \emph{inapplicable} rather than as 0\% accuracy. One set is testable; it is null (Appendix~\ref{app:external_audit}).}
\label{tab:external_audit}
\begin{scriptsize}
\setlength{\tabcolsep}{4pt}
\begin{tabular}{@{}p{1.28in}p{0.50in}p{0.60in}p{2.42in}@{}}
\toprule
\textbf{Rollout set} & \textbf{$n$} & \textbf{Fire rate} & \textbf{Criterion-4 verdict} \\
\midrule
Sandbagging, WMDP/MMLU & 1000 & 6.6--14.3\% & \textbf{Testable---and NULL}: $+4.1$\,pp ($p\!=\!0.32$), $+7.7$\,pp ($p\!=\!0.17$) \\
\addlinespace[1pt]
Sandbagging held-out tag; insider trading; paired-fact honesty; AI audit ($\times$2) & 366; 173; 40; 8, 16 & 0.0--12.5\% & Inapplicable: the rule is degenerate or underpowered in every cell---criterion 5 failing outside the format it was written for, not an accuracy of 0\% \\
\addlinespace[1pt]
Roleplaying ($\times$2); AI Liar & 371; 371; 54 & --- & \textbf{Ineligible}: 371/371/54 distinct system prompts---$E$ is not fixed, so intervention (A) is unavailable by construction \\
\addlinespace[1pt]
Goal-directed lying & 27 & --- & Ineligible: prompt materials, no model outputs \\
\bottomrule
\end{tabular}
\end{scriptsize}
\end{table}

Criterion 4 needs a corpus in which $D$ varies while $E$ does not, which our own materials cannot be. We audit ten rollout sets released with \citet{goldowskydill2025probes}, applying §\ref{sec:hedging_baseline}'s rule unchanged and offline and reading $D$ from each corpus's own per-rollout deception grade, which is distinct from the condition label. The sandbagging arm was inspected before pre-registration and is \textbf{exploratory}; a \textbf{pre-registered} arm on five held-out sets fixed the hypothesis, the pattern list, the text channel, both gates and all four reporting branches before any number was computed.

\paragraph{Result: one set can supply criterion 4, and it returns a null.} Three sets vary the system prompt per item and so cannot supply criterion 4 at all; five leave the rule inapplicable, firing on 0.0--12.5\% of trials with no cell both informative and of usable size---criterion 5 failing outside the format the rule was written for; one carries no model outputs. The single set clearing both gates is sandbagging ($n\!=\!1000$, one system prompt across all rollouts): at fixed condition label, realized deception moves the rule $+4.1$\,pp ($p\!=\!0.32$) and $+7.7$\,pp ($p\!=\!0.17$)---\textbf{null}. Restating $D$ without the corpus's LLM grader (final answer $\neq$ the item's key, 99.7\% parsed) is null as well, $+2.5$ and $+2.7$\,pp---criterion~4 needs an \emph{independently validated} proxy for $D$, not a second LLM judge reading the text the detector reads. \textbf{This is no evidence of a signal in the deployed-observable channel, not evidence that none exists}.

\paragraph{The channel a detector reads is itself an audit criterion.} On the \emph{full} transcript the same rule instead looks decisive---$+84.3$\,pp (92.0\% vs.\ 7.7\%, $p\!=\!0.0001$)---and the gap between those two numbers is the finding. The benchmark mandates a four-step scaffold whose system prompt declares steps 1--3 hidden from the user and asks for an ``\emph{incorrect} answer''; one pattern, \texttt{incorrect}, fires on 91.5\% of graded-deceptive full transcripts but 10.8\% of the same trials' visible answers, and no other exceeds 8\% anywhere. \textbf{The apparent effect is the model echoing its own instruction in a channel no deployed detector could read}, and it reproduces on the held-out tag ($+74.3$ vs.\ $+1.4$\,pp). That is §\ref{sec:identification}'s confound one level down---not the instruction producing the deception, but its \emph{wording} leaking into the scored text---and it adds a fourth requirement: a corpus supplies criterion 4 only if it (i)~fixes the elicitation prompt, (ii)~grades $D$ independently of the condition label, (iii)~observes the same item under both conditions, so incapability is separable from deception, and (iv)~exposes only the channel a deployed detector would read. \textbf{No audited set satisfies all four}: sandbagging meets (i)--(ii), fails (iii)--(iv), and every other set fails (i) or the applicability gate.
